Metastasis causes most cancer deaths, and the hybrid epithelial-mesenchymal state is strongly associated with metastatic competence and treatment resistance across carcinomas. How that state is measured in practice is a count: what fraction of cells occupy it when a sample is taken. This project asks whether the property that matters is instead a behavior, how a cell moves among states when conditions change, and whether two tumors with identical counts therefore carry different risk.
Answering that determines what is worth measuring. If the regimes differ in metastatic competence, a count is the wrong assay and one that perturbs a sample and observes the response is the right one. The project also identifies which molecular changes set one regime rather than another, which is what any attempt to move a tumor out of its most dangerous state would have to target.
The work does not deliver a clinical assay. It establishes whether the quantity is worth measuring, which is the question that must be answered before such an assay is justified, and it is designed so that a negative answer would be clear. Because the measurements build on breast cancer models and metastasis experiments already performed by our collaborators, the test can be run without waiting for new animal studies.